\section{Introduction}\label{sec:intro}

Problem 131 of Kagey's Open Problem Collection~\cite{kagey131} asks about a
Galton board where the ball tends to keep going the same way. In particular
it asks when the middle bin and the ends are equally likely. We count $n$
steps, so the bins are Kagey's row $n+1$, and we write $A_n$ for the number
of right steps. Unless we say otherwise, the first step is left or right
with probability $\frac12$ each. So the law of $A_n$ is symmetric about
$n/2$. For even $n$ put $h=n/2$, $C_n=\Prob(A_n=h)$ and
$E_n=\Prob(A_n=n)=\Prob(A_n=0)$.
Our companion paper~\cite{paperA}, which we call Paper~A, treats the Markov
walk. This walk repeats its last step with probability $p=1-q$. Paper~A
shows that there is exactly one $p_n$ with $C_n=E_n$ and that
\[
 n(1-p_n)=\log n-\tfrac12\log\log n+\tfrac12\log\frac{\pi}{8}+o(1).
\]
So a Markov walk needs about $\log n$ switches before the middle catches up
with the ends. We write $q_n=1-p_n$ and $z_n=nq_n$. A second companion
paper~\cite{paperB}, Paper~B, asks the same question for a Markov chain with
more than 2 directions.

The Markov walk remembers only its last step. We compare it with 2 walks
that have a longer memory. The elephant random walk of Sch\"utz and
Trimper~\cite{schutz2004} remembers its whole past. Each step copies a
uniformly chosen earlier step and flips it with probability $\varepsilon$.
In the second walk the steps come in runs of alternating direction, and the
runs after the first one have i.i.d.\ lengths $L$ with
$\Prob(L\ge k)=k^{-\alpha}$ for $k\ge1$. Its memory is the time since the
last switch. The first run is either a fresh run or, for $\alpha>1$, the
remaining part of a run in progress. We call these the fresh and the
stationary start. A remark and Appendix~\ref{app:aging} treat a third kind
of memory, aging, where the walk switches at step $k$ with probability
$c/k$.
Table~\ref{tab:compare} compares the answers.

\begin{table}[t]
\centering\small
\setlength{\tabcolsep}{4pt}
\begin{tabular}{@{}>{\raggedright\arraybackslash}p{1.65cm}
 >{\raggedright\arraybackslash}p{2.35cm}
 >{\raggedright\arraybackslash}p{4.9cm}
 >{\raggedright\arraybackslash}p{1.9cm}
 >{\raggedright\arraybackslash}p{4.3cm}@{}}
\toprule
Walk & Rule & Crossing & Center at the crossing & Shape at the crossing\\
\midrule
Markov \cite{paperA}
 & switch with probability $q$
 & $nq_n=\log n-\frac12\log\log n+\frac12\log\frac\pi8+o(1)$
 & order $(\log n)^{1/2}/n$
 & the middle and the ends are the largest bins, bins $1$ and $n-1$ the
 smallest\\
\addlinespace
Elephant (Sections \ref{sec:erw}, \ref{sec:shape})
 & copy a uniform earlier step, flip it with probability $\varepsilon$
 & $x_n=2\log n-\log\log n-{}$\newline
 $\log16+o(1)$, and\newline
 $x_n=2nq_n-\log(2\pi)+o(1)$
 & $\sim4x_n/n^2\sim8\log n/n^2$
 & each end below its neighbor, modes about $x_n\log x_n$ from the ends
 (for very large $n$, using one computed sign), the center below all
 interior bins outside a window (all of them if
 Conjecture~\ref{conj:convex} holds)\\
\addlinespace
Heavy-tailed runs (Section \ref{sec:heavy})
 & alternating runs with $\Prob(L\ge k)=k^{-\alpha}$
 & $R_n=E_n/C_n\sim C(\alpha)n^{1-\alpha+1/\alpha}$ for $1<\alpha<2$ and
 a stationary start, so the ends win for large $n$ exactly when
 $\alpha<\varphi=(1+\sqrt5)/2$. With a fresh start they win for $\alpha<1$ and lose for
 $\alpha>1$.
 & order $n^{-1/\alpha}$ for $1<\alpha<2$
 & a stable bulk of width $n^{1/\alpha}$ and 2 end atoms
 \cite{godrecheluck2001,zdk2015} (the full profile is not studied here)\\
\bottomrule
\end{tabular}
\caption{The 3 walks near the point where the middle bin and the ends
are equally likely. For the Markov and elephant walks, $n$ is the number of
steps and the crossing is a value of the switch or flip probability, with
$x_n=n\varepsilon_n$ for the elephant walk. For the heavy-tailed walk the
exponent $\alpha$ is fixed and we compare $E_n$ with $C_n$ as $n$ grows.}
\label{tab:compare}
\end{table}

\paragraph{Results.}
For the elephant walk, $C_n/E_n$ is strictly increasing in $\varepsilon$, so
the crossing $\varepsilon_n$ exists and is unique
(Proposition~\ref{prop:mlr}). The derivative of $C_n$ at $\varepsilon=0$ is
exactly $4/(n+2)$ (Proposition~\ref{prop:first-mutation}).
Theorem~\ref{thm:erw-bounds} gives explicit bounds for $C_n$. In particular
$\log C_n=\log\frac{4\varepsilon}{n+2}+2\varepsilon\log n+O(\varepsilon)$
when $\varepsilon(1+\log n)\le\frac1{100}$. Hence $x_n=n\varepsilon_n$
satisfies $x_n+\log x_n=2\log n-\log8+O((\log n)^2/n)$
(Theorem~\ref{thm:erw-crossing}). The reason for the factor 2 is that in
the elephant walk a single early flip already moves a large family of
steps. So the middle is reached with probability of order $\varepsilon/n$,
and the ends must fall to about $n^{-2}$ before they lose, instead of about
$n^{-1}$ for the Markov walk. Since $E_n\approx\frac12e^{-n\varepsilon}$,
the crossing is near $2\log n$ and not $\log n$. Using the third-order
expansion of Paper~A, Theorem~\ref{thm:erw-markov} shows that
$x_n-(2z_n-\log(2\pi))=c_*/\log n+O((\log\log n)^2/(\log n)^2)$ with
$c_*=\frac12\log(2\pi)-\frac14$. By Theorem~\ref{thm:K2} the second
coefficient of $C_n$ in $\varepsilon$ divided by the first is
$2\log n+K_2+o(1)$, where $K_2=2\gamma-2-\log2$ and $\gamma$ is Euler's
constant.

At the crossing and for large $n$, each end is below its neighbor
(Theorem~\ref{thm:dip} and Corollary~\ref{cor:dip}), so the law is not
unimodal. With $x=n\varepsilon$,
the number of steps that differ from the first step is close to a compound
Poisson law (Theorem~\ref{thm:cp}), with a Landau limit near $x\log x$ on
the scale $x$ (Theorem~\ref{thm:landau}). For very large $n$, and using one
numerically computed sign, the same theorem puts the modes near
$x\log x-0.2228x$. We also show $\Prob(A_n=n-m)\sim x/(2m^2)$ for
$x\log^2x\ll m\ll n$ (Theorem~\ref{thm:profile}). At the crossing and for
large $n$, the center is below every interior bin outside a window of width of order
$n^{3/4}(\log n)^{3/4}$ (Theorem~\ref{thm:shape}).
Conjecture~\ref{conj:convex} would close this window, and we checked it for
every $10\le n\le600$ and for 5 larger $n$.

For heavy-tailed runs with $1<\alpha<2$ and any first run,
Theorem~\ref{thm:center} gives $C_n$ to leading order, and it is of order
$n^{-1/\alpha}$. With a stationary start the end is of order $n^{1-\alpha}$,
so $R_n=E_n/C_n\sim C(\alpha)n^{1-\alpha+1/\alpha}$
(Corollary~\ref{cor:golden}). The reason is that the end needs one first run
that covers all $n$ steps, while the center needs the renewal processes of
$+$ and $-$ runs to meet, and these fluctuate on the scale $n^{1/\alpha}$.
The exponent vanishes at the golden ratio $\varphi=(1+\sqrt5)/2$, and there
$R_n\to C(\varphi)=0.7761317206\ldots$, so the center still wins.
Theorems~\ref{thm:alpha2} and~\ref{thm:lamperti} treat $\alpha\ge2$ and
$0<\alpha<1$. Remark~\ref{rem:crossover} gives computed crossover sizes for
$\alpha$ just below $\varphi$, and Remark~\ref{rem:heat} discusses
$\alpha=5/3$, the value used for anomalous heat conduction.
Remark~\ref{rem:aging} treats aging. The crossing is exactly $c=1$ when the
first step is fixed, and for a fair first step it is
$c_n=1-\log2/(\log n+\gamma+2-2\log2)+O((\log n)^{-3})$.

\paragraph{Prior work.}
The steps of the elephant walk split into clusters, since a step either
copies an earlier step exactly or starts a new cluster with a fair sign. This
is K\"ursten's coupling with bond percolation on random recursive
trees~\cite{kursten2016}, and Baur and Bertoin~\cite{baurbertoin2016} relate
the walk to P\'olya-type urns. Our regime $n\varepsilon\asymp\log n$ lies in
the strongly supercritical percolation regime of Baur~\cite{baur2016} (see
also Bertoin~\cite{bertoin2014}). Those results describe cluster sizes and
do not cover the center probability, which is a large deviation of the root
cluster. Gu\'erin, Laulin, Raschel and
Simon~\cite[proof of Theorem~1.1]{glrs2025} prove that the law with a fixed
first step is unimodal for every $n$ and every memory parameter. Their other
results fix the memory parameter as $n$ grows. We did not find an asymptotic
formula for the center probability in the literature when $\varepsilon\to0$
as $n$ grows.

For heavy-tailed runs, Godr\`eche and Luck~\cite{godrecheluck2001} give the
transform of the occupation time of an alternating renewal process and the
scalings of the center and of both edges. They also give Lamperti's law and
its critical value $\alpha_c$, which they credit to Baldassarri, Bouchaud,
Dornic and Godr\`eche. Zaburdaev, Denisov and Klafter~\cite{zdk2015} and
Dhar, Saito and Derrida~\cite{dsd2013} describe the profile of a L\'evy
walk. Wang, Li and H\"anggi~\cite{wlh2016} already observe that equal decay
of the central peak and the side peaks gives the golden ratio, as a
heuristic for the continuum model of heat conduction. Here we prove the
lattice statement with a local limit at the center, keep the 2 starts apart,
and compute $C(\alpha)$ and the crossover sizes. For aging, Dietz and
Sethuraman~\cite{ds2007} proved the limit law of $A_n/n$, Engl\"ander and
Volkov~\cite{ev2018} gave a second proof, and Engl\"ander, Volkov and
Wang~\cite{evw2020} found the scaling limit of the walk (see
Remark~\ref{rem:aging} and Appendix~\ref{app:aging}).

\paragraph{Organization.}
Sections~\ref{sec:erw} and~\ref{sec:shape} treat the crossing and the shape
of the elephant walk. Section~\ref{sec:heavy} treats heavy-tailed runs and
ends with the remark on aging. Section~\ref{sec:verify} describes the
computations and lists open questions. Appendices~\ref{app:erw},
\ref{app:shape} and~\ref{app:heavy} hold the longer proofs for
Sections~\ref{sec:erw} to~\ref{sec:heavy}, and Appendix~\ref{app:aging}
proves the aging crossing.
